\section{InternAgent Software Development and Application Interface}


\begin{figure}[t]
    \centering
    \includegraphics[width=\linewidth]{Figures/platform-internagent.pdf}
    \vspace{-4pt}
    \caption{InternAgent software platform includes features such as the user entry interface, task selection interface, idea-tree visualization and human-computer interaction interface, code generation, and auto-debug interface. In the near future, we plan to support additional functionalities, including custom dataset uploads and academic idea thinking modes.}
    \label{fig:showcase}
\end{figure}


\noindent \textbf{Software Development.} Fig.~\ref{fig:showcase} shows the front-end interface of the current InternAgent software platform. Overall, InternAgent software platform employs a frontend-backend separation design pattern, building a highly scalable distributed service platform. 

The frontend layer is developed based on the React framework, featuring an advanced visual interaction system. Key innovations include an infinite canvas rendering engine supporting multi-node topology, a collaborative mind mapping component driven by state synchronization, a code editor supporting multiple formats, and a real-time training metrics visualization dashboard.

The backend leverages a cloud-native technology stack, utilizing a dynamic container orchestration engine for elastic resource scheduling, a distributed asynchronous task queue for high concurrency support, and a cross-cloud storage gateway for data synchronization across heterogeneous cloud environments. Additionally, a microservice governance system is established using a Service Mesh architecture.

The entire system is delivered through containerization, with Kubernetes cluster management enabling self-healing and intelligent scaling, ensuring business continuity while significantly improving resource utilization efficiency.


\noindent \textbf{InternAgent Trial Application: Accelerating Scientific Research.} InternAgent is the first end-to-end Autonomous Scientific Research (ASR) multi-agent framework. The release of this innovative achievement enables Agent to autonomously analyze problems, review literature, conduct deep research, reflect, and carry out experiments, much like human scientists. It can automatically complete the entire process from hypothesis generation to experimental validation, starting from an initial research idea. This marks the official arrival of a new era in AI-driven scientific discovery.

Due to InternAgent's reliance on massive API key calls and GPU compute clusters for executing experiments, there are associated costs. As a result, we have recently launched a whitelist application page. On this page, applicants are required to provide information about the scenarios in which they will use InternAgent, their research tasks, and the institution they belong to, so that we can select the initial batch of seed users in InternAgent. The link to the application experience page is as follows: \texttt{\url{https://discovery.intern-ai.org.cn/}}.



